\section{规范投影、证书与结果归因}

\subsection{投影流水线}

\srcpath{project\_to\_canonical\_models} 的职责可分为：
\begin{enumerate}
  \item 聚合 charger、投影三相兼容视图并物化 DC storage；
  \item 把 2W/3W 变压器、开关、断路器和宏设备转为规范支路/设备；
  \item 展开能量路由器，并记录显式 expansion provenance；
  \item 按策略收缩理想或阈值小阻抗连接，组合 \texttt{BusMergeMap}；
  \item 剥离无源死岛、规范化 DC bus index，重写所有域限定引用；
  \item 建立支路展开、组件映射、投影报告和机器可读证书。
\end{enumerate}
具体顺序由 \srcpath{src/model/network\_utils.cpp:project\_to\_canonical\_models} 定义，
调用者不应自行重排其中步骤。

\subsection{投影策略与证书}

\texttt{ProjectionMode::ExactIdeal} 仅收缩作者阻抗严格为零或有明确 ideal provenance 的设备；
\texttt{ThresholdApproximate} 还允许收缩小于默认 $10^{-4}$ pu 阈值的非零设备。
\texttt{ProjectionCertificate} 记录 mode、threshold、每条 \texttt{MergeRecord} 的阻抗与
\texttt{ExactIdeal/ThresholdApproximate}、诊断、投影前 DC 母线数及死岛 position 映射。
\fld{exact()} 只表示没有 approximate merge，不代表宏设备投影或下游分析整体无近似。

tap 非 1 或 shift 非零的支路不允许按零阻抗普通连接收缩。\fld{ideal\_connectivity} 是作者意图证明，
用于区别理想拓扑边和数值很短的真实线路。

\subsection{BusMergeMap}

\texttt{BusMergeMap} 同时保存
\fld{ext\_to\_int}、\fld{ext\_to\_orig\_pos}、\fld{int\_to\_ext}、groups、dead bus、
原支路 position 到投影支路 position，以及需求/发电两套参与因子。
\begin{itemize}
  \item Intensive：把代表值广播到成员，死岛返回零；
  \item ExtensiveDemand：按 \fld{extensive\_participation} 拆分；
  \item ExtensiveGeneration：按 \fld{generation\_participation} 拆分；
  \item 映射声称已构建但关键表缺失时，反投影抛异常而不是猜测连续编号。
\end{itemize}
DC 死岛恢复使用独立 prestrip position 映射，不能借用 AC bus map。

\subsection{支路与组件 provenance}

\texttt{BranchExpandMap} 把规范 ACBranch 关联回 Transformer2W、Transformer3W pair、Switch 或
CircuitBreaker，并保留原端子和闭合状态。\texttt{ComponentMapping} 记录 source/canonical type、id 和
participation factor。稳定 ID 是字符串化来源；映射不是依赖名称前缀的隐式约定。

\subsection{可观测量恢复等级}

\texttt{ObservableAttribution} 按观测类型统计 canonical entity 和 attributed entity，并使用：
\begin{itemize}
  \item \texttt{Strong}：在声明模型下可确定恢复；
  \item \texttt{Approximate}：采用等分、近似合并或不可辨识约定；
  \item \texttt{AuditOnly}：可追溯但不是求解观测；
  \item \texttt{Unsupported}：没有可用恢复路径。
\end{itemize}
\fld{total()} 只要求 recovery 非 Unsupported 且计数相等；它不把 Approximate 提升为 Strong。

\subsection{设备终端流与符号}

\srcpath{src/model/result\_attribution.cpp:CanonicalToRichOperator::apply} 以 authored system、projected system、
PF/OPF 结果和 mapping 构建设备级结果。网络端子功率常采用“母线流向设备”符号，源荷署名常采用
“注入母线为正”符号；调用者必须读取结果类型的 terminal/sign contract，不能跨族直接相加。
停运组件统一归零；同母线多个理想开关/断路器不可辨识时按约定分配并降级恢复等级。

\subsection{归因索引与复杂度（AUD-097）}

\srcpath{src/model/result\_attribution.cpp} 在每次调用中建立 AC/DC 分离的电压索引、
稳定 ID 到规范容器 position 的索引，并在结果行生成后按实际查询的组件族惰性建立行索引。
不跨调用缓存状态，不改变结果行、来源记录顺序或恢复等级。
电压仍优先读取 authored position 对应的 PF 值，短向量退回作者 \fld{vm\_pu}，
不存在的 bus 保留原有零值回退；该回退不能作为已求解证据。
电气行匹配要求 type、domain 和 ID 一致；无 domain 字段的投影来源报告仍保留
同 type/ID 的全部匹配，这是既有 provenance 边界。

记母线数为 $B$、结果行数为 $C$、来源映射数为 $M$、发电机数为 $G$，
原有重复查找包含 $O(CB+MC+G^2)$ 工作量。
哈希索引使这些查找项变为期望 $O(B+FC+G+M+Q)$，其中 $F$ 为实际查询的固定组件族数，
$Q$ 包含查询和输出匹配数；端子流计算、剩余设备扫描和 OPF 模型复制不在此界内。
预设验收是 10000 节点归因耗时降低至少 70\%，且小系统无明显回退、全字段输出相同。
初版全量建索引使 14 节点增加约 $10\,\mu$s；将分配成本纳入模型后改为按族惰性建立，
最终 14 节点每次中位耗时为 $14.065\to13.748\,\mu$s。

\srcpath{tools/attribution\_performance\_benchmark.py} 对真实收敛 PF 后的固定状态计时，
每组运行三个新进程，小系统每进程平均 100 次归因。
10000 节点合成例为 $0.524050\to0.007305$ s，
\texttt{case9241pegase} 为 $0.184947\to0.005857$ s；全部序列化字段哈希一致。
完整 HTTP 响应另测为 $0.942449\to0.836189$ s，不能把局部加速倍数解释为整体加速。
\srcpath{tests/test\_result\_attribution.cpp} 覆盖乱序稀疏 ID、AC/DC 同号、
短/空/跨次变化 PF、规范 OPF position 与有序来源映射；4 用例、96 断言通过。
关联五目标共 99 用例、669 断言和 GUI API 82/82 通过。
归因与测试单元的 ASan/UBSan 通过 4 用例、96 断言；其他 Release 档案未插桩。
依赖 pin 不匹配且保留脏树，本轮仅增量编译链接，未声称全量构建或全库 sanitizer 通过。
详细命令、数值和剩余热点见 \srcpath{docs/testing/module\_code\_audit.md}。

\subsection{投影后审计不变量}

至少检查：AC/DC 引用均存在、AC/DC 同号不串域、映射 group 覆盖所有 surviving authored bus、
参与因子组内和为 1、死岛输出为零、支路展开来源可追、强度量往返一致、广延量总和守恒，以及
所有近似 merge 在证书中有记录。投影“成功返回”本身不替代这些不变量。
